\begin{beginnerstatement}[Kurz erklärt: Tests und Builds]

Ein automatisierter Test führt vorgegebene Prüfungen ohne manuelles
Durchklicken aus. Ein Unit-Test betrachtet einen kleinen, abgegrenzten Teil des
Codes. Ein Integrationstest prüft, ob mehrere Teile wie Backend und Datenbank
zusammenarbeiten. Ein E2E-Test verfolgt einen vollständigen Ablauf aus Sicht
eines Nutzers. Eine Testsuite bündelt zusammengehörige Tests; das Template
verwendet dafür Pytest bei Python, Vitest im Frontend, Playwright für E2E und
Cargo bei Rust. \emph{PASS} bedeutet bestanden, \emph{FAIL} fehlgeschlagen und
\emph{SKIP} übersprungen; ein übersprungener Test ist kein Erfolgsnachweis. Ein
Build verarbeitet den Quellcode zu einem auslieferbaren Artefakt, hier je nach
Pfad etwa zur Web-Ausgabe oder zu einem Desktop-Paket. Ein ZIP bündelt Dateien
in einem Archiv; ein Container-Image verpackt eine Laufzeitumgebung, aus der
ein Container gestartet werden kann.

\end{beginnerstatement}
